\documentclass[11pt,a4paper]{article}

% ------------------------------------------------
% Encodage / langue
% ------------------------------------------------
\usepackage[utf8]{inputenc}
\usepackage[T1]{fontenc}
\usepackage[french]{babel}

% ------------------------------------------------
% Mise en page
% ------------------------------------------------
\usepackage[a4paper,margin=2cm]{geometry}
\usepackage{lmodern}
\usepackage{microtype}
\usepackage{setspace}
\setstretch{1.08}

% ------------------------------------------------
% Maths
% ------------------------------------------------
\usepackage{amsmath,amssymb,amsthm,mathrsfs}
\usepackage{mathtools}
\usepackage{bm}

% ------------------------------------------------
% Couleurs / boîtes
% ------------------------------------------------
\usepackage{xcolor}
\usepackage[most]{tcolorbox}

% ------------------------------------------------
% Divers
% ------------------------------------------------
\usepackage{enumitem}
\usepackage{hyperref}
\usepackage{fancyhdr}

\hypersetup{
    colorlinks=true,
    linkcolor=blue!50!black,
    urlcolor=blue!50!black
}

\pagestyle{fancy}
\fancyhf{}
\fancyhead[L]{Corrigé détaillé --- Terminale spécialité mathématiques}
\fancyhead[R]{Raisonnement et logique}
\fancyfoot[C]{\thepage}

% ------------------------------------------------
% Commandes
% ------------------------------------------------
\newcommand{\R}{\mathbb{R}}
\newcommand{\C}{\mathbb{C}}
\newcommand{\N}{\mathbb{N}}

\definecolor{bleuclair}{RGB}{230,240,255}
\definecolor{bleufonce}{RGB}{20,70,140}
\definecolor{vertclair}{RGB}{235,250,235}
\definecolor{grisclair}{RGB}{245,245,245}

\newtcolorbox{chapitrebox}[1]{
    colback=bleuclair,
    colframe=bleufonce,
    boxrule=0.8pt,
    arc=2mm,
    title=\bfseries #1
}

\newtcolorbox{methodebox}{
    colback=vertclair,
    colframe=black!60,
    boxrule=0.5pt,
    arc=2mm,
    title=\bfseries Idée de méthode
}

\newtcolorbox{logiquebox}{
    colback=grisclair,
    colframe=black!60,
    boxrule=0.5pt,
    arc=2mm,
    title=\bfseries Point logique essentiel
}

\newtheorem*{remarque}{Remarque}

% ------------------------------------------------
% Document
% ------------------------------------------------
\begin{document}

\begin{center}
    {\LARGE \textbf{Corrigé détaillé}}\\[0.4em]
    {\large Banque d'exercices de Terminale}\\[0.3em]
    {\normalsize Avec mise en valeur des raisonnements, de la logique et des arguments clés}
\end{center}

\vspace{0.7em}

\begin{methodebox}
Dans tout ce corrigé, on cherche non seulement à \textbf{trouver la réponse}, mais aussi à \textbf{justifier pourquoi la méthode choisie est valable}.  
En particulier :
\begin{itemize}
    \item on annonce l'idée directrice ;
    \item on précise les hypothèses utiles ;
    \item on vérifie les conditions de validité ;
    \item on conclut proprement.
\end{itemize}
\end{methodebox}

% =========================================================
\begin{chapitrebox}{Chapitre 1 --- Fonctions}
\end{chapitrebox}

\subsection*{Exercice 1}
Soit \(f(x)=\dfrac{x^2-1}{x-1}\) définie sur \(\R\setminus\{1\}\).

\begin{methodebox}
On commence par factoriser le numérateur.  
Ensuite, il faut bien distinguer :
\begin{itemize}
    \item la \textbf{fonction simplifiée} ;
    \item la \textbf{fonction d'origine}, qui n'est pas définie en \(x=1\).
\end{itemize}
C'est une distinction de logique très importante.
\end{methodebox}

\textbf{1) Simplification}

On reconnaît une différence de deux carrés :
\[
x^2-1=(x-1)(x+1).
\]
Donc, pour tout \(x\neq 1\),
\[
f(x)=\frac{(x-1)(x+1)}{x-1}=x+1.
\]

\textbf{Conclusion :} sur \(\R\setminus\{1\}\), on a
\[
f(x)=x+1.
\]

\begin{logiquebox}
On n'a pas le droit d'écrire directement \(f(x)=x+1\) sur tout \(\R\), car la fonction initiale n'est \textbf{pas définie} en \(x=1\).
\end{logiquebox}

\textbf{2) Variations}

La fonction \(x\mapsto x+1\) est affine de coefficient directeur \(1\), donc strictement croissante sur \(\R\).  
Par conséquent, \(f\) est strictement croissante sur chacun des intervalles
\[
]-\infty,1[ \quad \text{et} \quad ]1,+\infty[.
\]

\textbf{3) Prolongement par continuité en \(x=1\)}

Comme pour \(x\neq 1\), \(f(x)=x+1\), on a
\[
\lim_{x\to 1} f(x)=\lim_{x\to 1}(x+1)=2.
\]
Donc on peut définir une nouvelle fonction \(g\) par
\[
g(x)=
\begin{cases}
\dfrac{x^2-1}{x-1} & \text{si } x\neq 1,\\
2 & \text{si } x=1.
\end{cases}
\]
Alors \(g\) est continue en \(x=1\).

\textbf{4) Interprétation graphique}

La courbe de \(f\) est celle de la droite d'équation
\[
y=x+1
\]
\textbf{privée du point} de coordonnées \((1,2)\).  
Le prolongement par continuité consiste à \textbf{reboucher le trou} en ajoutant ce point.

% ---------------------------------------------------------
\subsection*{Exercice 2}
Comparer \(e^\pi\) et \(\pi^e\).

\begin{methodebox}
On ne peut pas comparer directement deux nombres de formes différentes.  
L'idée classique est d'appliquer le logarithme, qui est strictement croissant sur \(]0,+\infty[\).
\end{methodebox}

Comme \(e^\pi>0\) et \(\pi^e>0\), comparer \(e^\pi\) et \(\pi^e\) revient à comparer leurs logarithmes :
\[
\ln(e^\pi)=\pi
\qquad \text{et} \qquad
\ln(\pi^e)=e\ln(\pi).
\]
Il faut donc comparer \(\pi\) et \(e\ln(\pi)\), ou encore
\[
\frac{\pi}{e} \quad \text{et} \quad \ln(\pi).
\]
Une méthode plus élégante consiste à étudier la fonction
\[
\varphi(x)=\frac{\ln x}{x}
\quad \text{sur } ]0,+\infty[.
\]

Calculons sa dérivée :
\[
\varphi'(x)=\frac{\frac1x\cdot x-\ln x}{x^2}
=\frac{1-\ln x}{x^2}.
\]
Comme \(x^2>0\), le signe de \(\varphi'(x)\) est celui de \(1-\ln x\). Donc :
\[
\varphi'(x)>0 \iff \ln x<1 \iff x<e,
\]
\[
\varphi'(x)=0 \iff x=e,
\]
\[
\varphi'(x)<0 \iff x>e.
\]
Ainsi, \(\varphi\) est croissante sur \(]0,e[\), puis décroissante sur \(]e,+\infty[\).  
Elle atteint donc son maximum en \(x=e\).

Or \(\pi>e\), donc, comme \(\varphi\) décroît après \(e\),
\[
\varphi(\pi)<\varphi(e).
\]
Ainsi,
\[
\frac{\ln \pi}{\pi}<\frac{\ln e}{e}=\frac{1}{e}.
\]
En multipliant par \(\pi e>0\), on obtient
\[
e\ln \pi<\pi.
\]
En exponentiant,
\[
\pi^e<e^\pi.
\]

\textbf{Conclusion :}
\[
\boxed{e^\pi>\pi^e.}
\]

% ---------------------------------------------------------
\subsection*{Exercice 3}
Soit \(f(x)=\dfrac{x}{x^2+1}\).

\textbf{1) Variations}

La fonction est définie sur \(\R\), car \(x^2+1>0\) pour tout réel \(x\).

Calculons la dérivée :
\[
f'(x)=\frac{(x^2+1)\cdot 1-x\cdot 2x}{(x^2+1)^2}
=\frac{x^2+1-2x^2}{(x^2+1)^2}
=\frac{1-x^2}{(x^2+1)^2}.
\]
Le dénominateur est toujours strictement positif.  
Le signe de \(f'(x)\) est donc celui de \(1-x^2=(1-x)(1+x)\).

Ainsi :
\[
f'(x)>0 \text{ sur } ]-1,1[,\qquad
f'(x)=0 \text{ pour } x=\pm1,\qquad
f'(x)<0 \text{ sur } ]-\infty,-1[ \cup ]1,+\infty[.
\]

Donc :
\begin{itemize}
    \item \(f\) décroît sur \(]-\infty,-1]\),
    \item \(f\) croît sur \([-1,1]\),
    \item \(f\) décroît sur \([1,+\infty[\).
\end{itemize}

Calculons les valeurs remarquables :
\[
f(-1)=\frac{-1}{2}=-\frac12,\qquad f(1)=\frac12.
\]

\textbf{2) Encadrement}

Les variations montrent que la valeur minimale est \(-\dfrac12\) et la valeur maximale est \(\dfrac12\).  
Donc
\[
\boxed{-\frac12\leq \frac{x}{x^2+1}\leq \frac12.}
\]

\textbf{3) Sans dérivée}

On veut montrer
\[
\frac{x}{x^2+1}\leq \frac12.
\]
Comme \(x^2+1>0\), cela équivaut à
\[
2x\leq x^2+1
\iff x^2-2x+1\geq 0
\iff (x-1)^2\geq 0,
\]
ce qui est vrai.

De même,
\[
\frac{x}{x^2+1}\geq -\frac12
\]
équivaut à
\[
2x\geq -(x^2+1)
\iff x^2+2x+1\geq 0
\iff (x+1)^2\geq 0,
\]
ce qui est vrai.

Donc
\[
-\frac12\leq \frac{x}{x^2+1}\leq \frac12.
\]

\begin{logiquebox}
Cette deuxième méthode est très intéressante : elle montre qu'une inégalité peut parfois se démontrer \textbf{plus élégamment par algèbre} que par étude de fonction.
\end{logiquebox}

% ---------------------------------------------------------
\subsection*{Exercice 4}
Soit \(f(x)=\sqrt{x+2}-\sqrt{x}\) sur \([0,+\infty[\).

\textbf{1) Variations}

Pour \(x\geq 0\), la fonction est bien définie.  
Calculons sa dérivée :
\[
f'(x)=\frac{1}{2\sqrt{x+2}}-\frac{1}{2\sqrt{x}}.
\]
Pour \(x>0\), comme \(\sqrt{x+2}>\sqrt{x}\), on a
\[
\frac{1}{2\sqrt{x+2}}<\frac{1}{2\sqrt{x}},
\]
donc
\[
f'(x)<0.
\]
Ainsi, \(f\) est strictement décroissante sur \(]0,+\infty[\), donc sur \([0,+\infty[\).

\textbf{2) Limite en \(+\infty\)}

On rationalise :
\[
f(x)=\sqrt{x+2}-\sqrt{x}
=\frac{(x+2)-x}{\sqrt{x+2}+\sqrt{x}}
=\frac{2}{\sqrt{x+2}+\sqrt{x}}.
\]
Quand \(x\to +\infty\), le dénominateur tend vers \(+\infty\), donc
\[
\lim_{x\to+\infty} f(x)=0.
\]

\textbf{3) Positivité}

Pour tout \(x\geq 0\), on a \(x+2>x\), donc
\[
\sqrt{x+2}>\sqrt{x}.
\]
Par conséquent,
\[
f(x)>0.
\]

\textbf{4) Plus grand réel \(M\)}

Comme \(f\) est décroissante sur \([0,+\infty[\), elle atteint sa valeur maximale en \(x=0\) :
\[
f(0)=\sqrt{2}-0=\sqrt2.
\]
Ainsi, le plus grand réel \(M\) tel que \(f(x)\leq M\) pour tout \(x\geq 0\) est
\[
\boxed{M=\sqrt2.}
\]

% ---------------------------------------------------------
\subsection*{Exercice 5}
Résoudre \(x+\sqrt{x}=2\).

\begin{methodebox}
Dès qu'apparaît \(\sqrt{x}\), il faut penser au domaine de définition : ici \(x\geq 0\).
\end{methodebox}

\textbf{1) Méthode algébrique}

Posons
\[
t=\sqrt{x}\geq 0.
\]
Alors \(x=t^2\), et l'équation devient
\[
t^2+t=2
\iff t^2+t-2=0
\iff (t+2)(t-1)=0.
\]
Comme \(t\geq 0\), on ne garde que
\[
t=1.
\]
Donc
\[
x=t^2=1.
\]

\textbf{2) Par étude de fonction}

On considère
\[
g(x)=x+\sqrt{x}
\quad \text{sur } [0,+\infty[.
\]
La fonction \(x\mapsto x\) est croissante, la fonction \(x\mapsto \sqrt{x}\) aussi, donc \(g\) est strictement croissante sur \([0,+\infty[\).  
Ainsi, l'équation \(g(x)=2\) admet au plus une solution.

Or
\[
g(1)=1+1=2.
\]
Donc \(x=1\) est l'unique solution.

\textbf{3) Intérêt des deux méthodes}

La méthode algébrique est rapide.  
La méthode par étude de fonction donne une information plus qualitative : elle explique \textbf{pourquoi} la solution est unique.

\[
\boxed{x=1}
\]

% =========================================================
\begin{chapitrebox}{Chapitre 2 --- Suites}
\end{chapitrebox}

\subsection*{Exercice 1}
\[
u_0=1,\qquad u_{n+1}=\frac{2u_n+3}{u_n+4}.
\]

\textbf{1) Suite bien définie}

Il faut vérifier que le dénominateur \(u_n+4\) n'est jamais nul.

Montrons d'abord que \(u_n>-4\) pour tout \(n\).  
On observe déjà que \(u_0=1>-4\).

Supposons \(u_n>-4\). Alors \(u_n+4>0\), donc \(u_{n+1}\) est bien défini.  
De plus,
\[
u_{n+1}+4=\frac{2u_n+3}{u_n+4}+4
=\frac{2u_n+3+4u_n+16}{u_n+4}
=\frac{6u_n+19}{u_n+4}.
\]
Si \(u_n>-4\), alors \(u_n+4>0\), et
\[
6u_n+19>6(-4)+19=-5.
\]
Cela ne suffit pas directement à conclure que \(u_{n+1}>-4\).  
Essayons plutôt de montrer une propriété plus forte : \(u_n>0\).

On a \(u_0=1>0\).  
Si \(u_n>0\), alors
\[
2u_n+3>0,\qquad u_n+4>0,
\]
donc
\[
u_{n+1}=\frac{2u_n+3}{u_n+4}>0.
\]
Par récurrence,
\[
u_n>0 \quad \text{pour tout } n.
\]
En particulier, \(u_n+4>0\), donc la suite est bien définie.

\textbf{2) Conjecture de limite}

Calculons :
\[
u_1=\frac{2\cdot1+3}{1+4}=\frac55=1.
\]
Donc ici,
\[
u_1=1=u_0.
\]
On conjecture que la suite est constante égale à \(1\).

\textbf{3) Si \(u_n\to \ell\), équation vérifiée}

Si \((u_n)\) converge vers \(\ell\), alors en passant à la limite dans la relation de récurrence,
\[
\ell=\frac{2\ell+3}{\ell+4}.
\]
On résout :
\[
\ell(\ell+4)=2\ell+3
\iff \ell^2+4\ell=2\ell+3
\iff \ell^2+2\ell-3=0
\iff (\ell-1)(\ell+3)=0.
\]
Donc les limites possibles sont
\[
\ell=1 \quad \text{ou} \quad \ell=-3.
\]

\textbf{4) Convergence}

Montrons que la suite est constante.  
Supposons \(u_n=1\). Alors
\[
u_{n+1}=\frac{2\cdot1+3}{1+4}=1.
\]
Comme \(u_0=1\), la récurrence donne
\[
u_n=1 \quad \text{pour tout } n.
\]
Donc la suite converge vers \(1\).

\[
\boxed{u_n=1 \text{ pour tout } n,\quad \lim u_n=1.}
\]

% ---------------------------------------------------------
\subsection*{Exercice 2}
\[
u_0\in ]0,1[,\qquad u_{n+1}=u_n(2-u_n).
\]

\textbf{1) Montrer que \(0<u_n<1\)}

On procède par récurrence.

\emph{Initialisation :} c'est vrai par hypothèse pour \(u_0\).

\emph{Hérédité :} supposons \(0<u_n<1\). Alors
\[
1<2-u_n<2.
\]
Comme \(u_n>0\), on obtient
\[
u_{n+1}=u_n(2-u_n)>0.
\]
Et comme \(2-u_n<2\),
\[
u_{n+1}<2u_n.
\]
Cela ne suffit pas pour montrer \(u_{n+1}<1\).  
On écrit plutôt
\[
u_{n+1}=u_n(2-u_n)=2u_n-u_n^2.
\]
Alors
\[
u_{n+1}-1=2u_n-u_n^2-1=-(u_n-1)^2\leq 0.
\]
Comme \(u_n\neq 1\), on a même
\[
u_{n+1}<1.
\]
Donc
\[
0<u_{n+1}<1.
\]
La propriété est héréditaire.

Ainsi,
\[
0<u_n<1 \quad \forall n\in\N.
\]

\textbf{2) Sens de variation}

Calculons
\[
u_{n+1}-u_n=u_n(2-u_n)-u_n=u_n(1-u_n).
\]
Or \(0<u_n<1\), donc
\[
u_n(1-u_n)>0.
\]
Ainsi,
\[
u_{n+1}-u_n>0.
\]
Donc la suite est strictement croissante.

\textbf{3) Limite}

La suite est croissante et majorée par \(1\), donc elle converge vers une limite \(\ell\in ]0,1]\).

En passant à la limite dans la relation,
\[
\ell=\ell(2-\ell).
\]
Donc
\[
\ell=2\ell-\ell^2
\iff \ell-\ell^2=0
\iff \ell(1-\ell)=0.
\]
Donc \(\ell=0\) ou \(\ell=1\).  
Mais comme \(u_n>0\) et \((u_n)\) croît, la limite ne peut pas être \(0\). Donc
\[
\boxed{\lim u_n=1.}
\]

% ---------------------------------------------------------
\subsection*{Exercice 3}
\[
u_n=\sum_{k=1}^n \frac{1}{k(k+1)}.
\]

\textbf{1) Premiers termes}

\[
u_1=\frac{1}{1\cdot2}=\frac12,
\]
\[
u_2=\frac12+\frac{1}{2\cdot3}=\frac12+\frac16=\frac23.
\]

\textbf{2) Simplification}

On décompose :
\[
\frac{1}{k(k+1)}=\frac{1}{k}-\frac{1}{k+1}.
\]
En effet,
\[
\frac{1}{k}-\frac{1}{k+1}
=\frac{k+1-k}{k(k+1)}
=\frac{1}{k(k+1)}.
\]
Donc
\[
u_n=\sum_{k=1}^n \left(\frac{1}{k}-\frac{1}{k+1}\right).
\]
La somme est télescopique :
\[
u_n=\left(1-\frac12\right)+\left(\frac12-\frac13\right)+\cdots+\left(\frac1n-\frac{1}{n+1}\right).
\]
Tous les termes intermédiaires se simplifient, il reste :
\[
u_n=1-\frac{1}{n+1}=\frac{n}{n+1}.
\]

\textbf{3) Limite}

\[
u_n=\frac{n}{n+1}=\frac{1}{1+\frac1n}\longrightarrow 1.
\]

\[
\boxed{u_n=\frac{n}{n+1}\quad \text{et} \quad \lim u_n=1.}
\]

% ---------------------------------------------------------
\subsection*{Exercice 4}
\[
u_0=1,\qquad u_{n+1}=\sqrt{2+u_n}.
\]

\textbf{1) Montrer que \(u_n\leq 2\)}

Par récurrence.

Initialement,
\[
u_0=1\leq 2.
\]
Supposons \(u_n\leq 2\). Alors
\[
2+u_n\leq 4
\]
et comme la racine carrée est croissante,
\[
u_{n+1}=\sqrt{2+u_n}\leq \sqrt4=2.
\]
Donc \(u_n\leq 2\) pour tout \(n\).

\textbf{2) Monotonie}

On veut comparer \(u_{n+1}\) et \(u_n\).  
Comme les termes sont positifs, on peut comparer leurs carrés :
\[
u_{n+1}\geq u_n
\iff \sqrt{2+u_n}\geq u_n
\iff 2+u_n\geq u_n^2.
\]
Cela équivaut à
\[
u_n^2-u_n-2\leq 0
\iff (u_n-2)(u_n+1)\leq 0.
\]
Or \(u_n\leq 2\) et \(u_n+1>0\), donc cette inégalité est vraie. Ainsi
\[
u_{n+1}\geq u_n.
\]
La suite est croissante.

\textbf{3) Limite}

La suite est croissante et majorée par \(2\), donc elle converge vers une limite \(\ell\leq 2\).

En passant à la limite :
\[
\ell=\sqrt{2+\ell}.
\]
Comme \(\ell\geq 0\), on peut élever au carré :
\[
\ell^2=2+\ell
\iff \ell^2-\ell-2=0
\iff (\ell-2)(\ell+1)=0.
\]
Les solutions sont \(2\) et \(-1\). Comme \(\ell\geq 0\), on garde
\[
\boxed{\ell=2.}
\]

% ---------------------------------------------------------
\subsection*{Exercice 5}
\[
u_n=\frac{n}{2^n}.
\]

\textbf{1) Sens de variation}

Comparons deux termes consécutifs :
\[
\frac{u_{n+1}}{u_n}
=\frac{\frac{n+1}{2^{n+1}}}{\frac{n}{2^n}}
=\frac{n+1}{2n}.
\]
Donc
\[
u_{n+1}\leq u_n \iff \frac{n+1}{2n}\leq 1
\iff n+1\leq 2n
\iff 1\leq n.
\]
Ainsi, la suite est décroissante à partir de \(n=1\).

\textbf{2) Limite}

On sait que l'exponentielle \(2^n\) croît beaucoup plus vite que \(n\). Donc
\[
\lim_{n\to+\infty}\frac{n}{2^n}=0.
\]
On peut aussi le voir en remarquant que la suite est positive, décroissante à partir d'un certain rang, donc convergente, et que sa limite ne peut être que \(0\).

\textbf{3) Existence de \(N\)}

On cherche \(n\) tel que
\[
\frac{n}{2^n}\leq 10^{-3}.
\]
Par essais :
\[
u_{10}=\frac{10}{1024}\approx 9{,}8\times 10^{-3},
\]
\[
u_{14}=\frac{14}{16384}\approx 8{,}5\times 10^{-4}<10^{-3}.
\]
Donc on peut prendre par exemple
\[
\boxed{N=14.}
\]

% =========================================================
\begin{chapitrebox}{Chapitre 3 --- Dérivation et convexité}
\end{chapitrebox}

\subsection*{Exercice 1}
\[
f(x)=x^3-3x.
\]

\textbf{1) Variations}

\[
f'(x)=3x^2-3=3(x^2-1)=3(x-1)(x+1).
\]
Donc :
\begin{itemize}
    \item \(f'(x)>0\) si \(x<-1\) ou \(x>1\),
    \item \(f'(x)<0\) si \(-1<x<1\).
\end{itemize}
Ainsi :
\begin{itemize}
    \item \(f\) croît sur \(]-\infty,-1]\),
    \item \(f\) décroît sur \([-1,1]\),
    \item \(f\) croît sur \([1,+\infty[\).
\end{itemize}

Valeurs remarquables :
\[
f(-1)=-1+3=2,\qquad f(1)=1-3=-2.
\]

\textbf{2) Convexité}

\[
f''(x)=6x.
\]
Donc :
\begin{itemize}
    \item \(f''(x)<0\) si \(x<0\) : la courbe est concave ;
    \item \(f''(x)>0\) si \(x>0\) : la courbe est convexe.
\end{itemize}

\textbf{3) Tangentes horizontales}

Une tangente est horizontale lorsque \(f'(x)=0\), donc pour
\[
x=-1 \quad \text{ou} \quad x=1.
\]
Les points correspondants sont \((-1,2)\) et \((1,-2)\).

\textbf{4) Point d'inflexion}

Il y a changement de signe de \(f''\) en \(x=0\), donc \(x=0\) est l'abscisse d'un point d'inflexion.  
Comme
\[
f(0)=0,
\]
le point d'inflexion est
\[
\boxed{(0,0).}
\]

% ---------------------------------------------------------
\subsection*{Exercice 2}
\[
f(x)=\ln x \quad \text{sur } ]0,+\infty[.
\]

\textbf{1) Courbe sous ses tangentes}

On a
\[
f'(x)=\frac1x,\qquad f''(x)=-\frac1{x^2}<0.
\]
Donc \(f\) est concave sur \(]0,+\infty[\).

Or une fonction concave est située \textbf{en dessous de chacune de ses tangentes}.  
C'est précisément une propriété fondamentale de concavité.

\textbf{2) Déduction de \(\ln x\leq x-1\)}

Considérons la tangente en \(x=1\).  
On a
\[
f(1)=\ln 1=0,\qquad f'(1)=1.
\]
L'équation de la tangente en \(1\) est donc
\[
y=f(1)+f'(1)(x-1)=x-1.
\]
Comme la courbe de \(\ln x\) est en dessous de cette tangente,
\[
\boxed{\ln x\leq x-1 \quad \text{pour tout } x>0.}
\]

% ---------------------------------------------------------
\subsection*{Exercice 3}
\[
f(x)=e^x.
\]

\textbf{1) Tangente en \(0\)}

\[
f(0)=1,\qquad f'(x)=e^x,\qquad f'(0)=1.
\]
L'équation de la tangente en \(0\) est
\[
y=1+x.
\]

\textbf{2) Comparaison entre \(e^x\) et \(x+1\)}

Posons
\[
g(x)=e^x-(x+1).
\]
Alors
\[
g'(x)=e^x-1.
\]
On a :
\begin{itemize}
    \item \(g'(x)<0\) si \(x<0\),
    \item \(g'(0)=0\),
    \item \(g'(x)>0\) si \(x>0\).
\end{itemize}
Donc \(g\) admet un minimum en \(x=0\). Or
\[
g(0)=1-(0+1)=0.
\]
Ainsi,
\[
g(x)\geq 0 \quad \forall x\in\R.
\]
Donc
\[
\boxed{e^x\geq x+1 \quad \forall x\in\R.}
\]

\textbf{3) Signe de \(e^x-x-1\)}

C'est exactement \(g(x)\), donc
\[
\boxed{e^x-x-1\geq 0,}
\]
avec égalité seulement pour \(x=0\).

% ---------------------------------------------------------
\subsection*{Exercice 4}
\[
f(x)=\frac1x \quad \text{sur } ]0,+\infty[.
\]

\textbf{1) Convexité}

\[
f'(x)=-\frac1{x^2},\qquad f''(x)=\frac{2}{x^3}>0 \quad (x>0).
\]
Donc \(f\) est convexe sur \(]0,+\infty[\).

\textbf{2) Inégalité}

Pour une fonction convexe,
\[
f\left(\frac{a+b}{2}\right)\leq \frac{f(a)+f(b)}{2}.
\]
En appliquant cette propriété à \(f(x)=1/x\), on obtient :
\[
\frac{1}{\frac{a+b}{2}} \leq \frac12\left(\frac1a+\frac1b\right).
\]

\[
\boxed{\frac{1}{\frac{a+b}{2}} \leq \frac12\left(\frac1a+\frac1b\right).}
\]

\begin{remarque}
C'est un cas particulier d'inégalité entre moyenne harmonique et moyenne arithmétique.
\end{remarque}

% ---------------------------------------------------------
\subsection*{Exercice 5}
\[
f(x)=x\ln x \quad \text{sur } ]0,+\infty[.
\]

\textbf{1) Variations}

\[
f'(x)=\ln x+1.
\]
Donc
\[
f'(x)=0 \iff \ln x=-1 \iff x=e^{-1}.
\]
Ainsi :
\begin{itemize}
    \item \(f'(x)<0\) si \(0<x<e^{-1}\),
    \item \(f'(x)>0\) si \(x>e^{-1}\).
\end{itemize}
Donc \(f\) décroît sur \(]0,e^{-1}]\), puis croît sur \([e^{-1},+\infty[\).

\textbf{2) Convexité}

\[
f''(x)=\frac1x>0 \quad (x>0).
\]
Donc \(f\) est convexe sur \(]0,+\infty[\).

\textbf{3) Unicité de solution de \(x\ln x=1\) sur \(]1,+\infty[\)}

Sur \(]1,+\infty[\), on a \(\ln x>0\), donc \(f'(x)=\ln x+1>0\).  
Ainsi, \(f\) est strictement croissante sur \(]1,+\infty[\).

Or
\[
f(1)=1\cdot \ln 1=0<1,
\]
et
\[
f(e)=e\cdot 1=e>1.
\]
Par continuité et croissance stricte, l'équation
\[
x\ln x=1
\]
admet une unique solution sur \(]1,+\infty[\).

% =========================================================
\begin{chapitrebox}{Chapitre 4 --- Logarithme et exponentielle}
\end{chapitrebox}

\subsection*{Exercice 1}
\[
e^{2x}-3e^x+2=0.
\]

Posons
\[
X=e^x.
\]
Comme \(e^x>0\), on a \(X>0\). L'équation devient
\[
X^2-3X+2=0
\iff (X-1)(X-2)=0.
\]
Donc
\[
X=1 \quad \text{ou} \quad X=2.
\]
Ainsi
\[
e^x=1 \Rightarrow x=0,
\qquad
e^x=2 \Rightarrow x=\ln 2.
\]

\[
\boxed{x=0 \quad \text{ou} \quad x=\ln 2.}
\]

% ---------------------------------------------------------
\subsection*{Exercice 2}
\[
\ln(x+1)+\ln(x-2)=\ln 4.
\]

\textbf{1) Domaine}

Il faut avoir
\[
x+1>0 \quad \text{et} \quad x-2>0,
\]
soit
\[
x>-1 \quad \text{et} \quad x>2.
\]
Donc l'ensemble de définition est
\[
]2,+\infty[.
\]

\textbf{2) Résolution}

Sur ce domaine, on peut additionner les logarithmes :
\[
\ln\big((x+1)(x-2)\big)=\ln 4.
\]
Comme \(\ln\) est injective sur \(]0,+\infty[\),
\[
(x+1)(x-2)=4.
\]
Développons :
\[
x^2-x-2=4
\iff x^2-x-6=0
\iff (x-3)(x+2)=0.
\]
Les solutions algébriques sont
\[
x=3 \quad \text{ou} \quad x=-2.
\]

\textbf{3) Vérification}

Or le domaine est \(]2,+\infty[\), donc seule la solution
\[
x=3
\]
est admissible.

\[
\boxed{x=3.}
\]

\begin{logiquebox}
En présence de logarithmes, il faut toujours \textbf{vérifier les conditions d'existence} avant et après la résolution.
\end{logiquebox}

% ---------------------------------------------------------
\subsection*{Exercice 3}
Étudier selon \(m\in\R\) le nombre de solutions de
\[
\ln x=mx.
\]

On travaille sur \(x>0\).  
L'équation équivaut à
\[
\frac{\ln x}{x}=m.
\]
On retrouve la fonction
\[
\varphi(x)=\frac{\ln x}{x}.
\]
On a déjà vu que \(\varphi\) admet un maximum en \(x=e\), égal à
\[
\varphi(e)=\frac1e.
\]
De plus,
\[
\lim_{x\to 0^+}\varphi(x)=-\infty,\qquad
\lim_{x\to+\infty}\varphi(x)=0.
\]
Donc l'image de \(\varphi\) est
\[
]-\infty,\frac1e].
\]

Ainsi :
\begin{itemize}
    \item si \(m>\dfrac1e\), aucune solution ;
    \item si \(m=\dfrac1e\), une unique solution \(x=e\) ;
    \item si \(0<m<\dfrac1e\), deux solutions ;
    \item si \(m\leq 0\), une unique solution.
\end{itemize}

\textbf{Justification du dernier cas :}  
sur \(]0,1[\), \(\ln x<0\), donc \(\varphi(x)<0\) ; de plus \(\varphi\) y est croissante de \(-\infty\) vers \(0\), donc chaque valeur négative est prise une seule fois. Pour \(m=0\), la solution est \(x=1\).

% ---------------------------------------------------------
\subsection*{Exercice 4}
Étudier
\[
f(x)=\frac{\ln x}{x}.
\]

On a
\[
f'(x)=\frac{1-\ln x}{x^2}.
\]
Le dénominateur étant strictement positif, le signe de \(f'\) est celui de \(1-\ln x\).  
Donc :
\begin{itemize}
    \item \(f\) croît sur \(]0,e]\),
    \item \(f\) décroît sur \([e,+\infty[\).
\end{itemize}
Le maximum est
\[
f(e)=\frac1e.
\]

Pour comparer \(e^\pi\) et \(\pi^e\), on observe que \(\pi>e\).  
Comme \(f\) décroît sur \([e,+\infty[\),
\[
f(\pi)<f(e),
\]
soit
\[
\frac{\ln\pi}{\pi}<\frac1e.
\]
D'où
\[
e\ln\pi<\pi.
\]
En exponentiant,
\[
\pi^e<e^\pi.
\]

\[
\boxed{e^\pi>\pi^e.}
\]

% ---------------------------------------------------------
\subsection*{Exercice 5}
Montrer que
\[
\ln x\leq x-1 \quad (x>0).
\]

Considérons
\[
g(x)=x-1-\ln x.
\]
Alors
\[
g'(x)=1-\frac1x=\frac{x-1}{x}.
\]
Donc :
\begin{itemize}
    \item \(g'(x)<0\) sur \(]0,1[\),
    \item \(g'(1)=0\),
    \item \(g'(x)>0\) sur \(]1,+\infty[\).
\end{itemize}
Ainsi \(g\) admet un minimum en \(x=1\). Or
\[
g(1)=1-1-\ln 1=0.
\]
Donc
\[
g(x)\geq 0 \quad \forall x>0,
\]
c'est-à-dire
\[
\boxed{\ln x\leq x-1.}
\]
L'égalité a lieu si et seulement si \(x=1\).

% =========================================================
\begin{chapitrebox}{Chapitre 5 --- Intégration et primitives}
\end{chapitrebox}

\subsection*{Exercice 1}
\[
\int_0^1 (2x+3)\,dx.
\]

Une primitive de \(2x+3\) est
\[
F(x)=x^2+3x.
\]
Donc
\[
\int_0^1 (2x+3)\,dx=F(1)-F(0)=(1+3)-0=4.
\]

\textbf{Interprétation géométrique :}  
La courbe \(y=2x+3\) est une droite. Sur \([0,1]\), l'aire sous la courbe est celle d'un trapèze de bases \(3\) et \(5\), de hauteur \(1\) :
\[
\mathcal{A}=\frac{3+5}{2}\times 1=4.
\]

\[
\boxed{\int_0^1 (2x+3)\,dx=4.}
\]

% ---------------------------------------------------------
\subsection*{Exercice 2}
\[
\int_0^1 xe^x\,dx.
\]

On peut chercher une primitive de \(xe^x\).  
On sait que
\[
(x-1)e^x
\]
est une primitive, car
\[
\frac{d}{dx}\big((x-1)e^x\big)=e^x+(x-1)e^x=xe^x.
\]
Ainsi
\[
\int_0^1 xe^x\,dx=\left[(x-1)e^x\right]_0^1.
\]
On calcule :
\[
(1-1)e^1=0,\qquad (0-1)e^0=-1.
\]
Donc
\[
\int_0^1 xe^x\,dx=0-(-1)=1.
\]

\[
\boxed{\int_0^1 xe^x\,dx=1.}
\]

% ---------------------------------------------------------
\subsection*{Exercice 3}
\[
f(x)=\frac1{x^2+1}.
\]

\textbf{1) Encadrement de \(f\) sur \([0,1]\)}

Si \(x\in[0,1]\), alors
\[
0\leq x^2\leq 1,
\]
donc
\[
1\leq x^2+1\leq 2.
\]
Comme l'inverse renverse l'ordre pour des nombres positifs,
\[
\frac12\leq \frac1{x^2+1}\leq 1.
\]

\textbf{2) Encadrement de l'intégrale}

En intégrant membre à membre sur \([0,1]\),
\[
\int_0^1 \frac12\,dx \leq \int_0^1 \frac1{x^2+1}\,dx \leq \int_0^1 1\,dx.
\]
Donc
\[
\frac12 \leq \int_0^1 \frac1{x^2+1}\,dx \leq 1.
\]

\[
\boxed{\frac12 \leq \int_0^1 \frac1{x^2+1}\,dx \leq 1.}
\]

% ---------------------------------------------------------
\subsection*{Exercice 4}
\[
A(x)=\int_0^x (t^2-2t+3)\,dt.
\]

\textbf{1) Expression de \(A(x)\)}

Une primitive de \(t^2-2t+3\) est
\[
\frac{t^3}{3}-t^2+3t.
\]
Donc
\[
A(x)=\left[\frac{t^3}{3}-t^2+3t\right]_0^x
=\frac{x^3}{3}-x^2+3x.
\]

\textbf{2) Variations}

\[
A'(x)=x^2-2x+3=(x-1)^2+2>0 \quad \forall x\in\R.
\]
Donc \(A\) est strictement croissante sur \(\R\).

\textbf{3) Interprétation}

\(A(x)\) représente l'aire algébrique comprise entre la courbe
\[
y=t^2-2t+3
\]
et l'axe des abscisses, entre \(0\) et \(x\).  
Comme l'intégrande vaut toujours
\[
(t-1)^2+2>0,
\]
cette aire algébrique est en fait une aire géométrique positive.

% ---------------------------------------------------------
\subsection*{Exercice 5}
\[
I_n=\int_0^1 x^n\,dx.
\]

\textbf{1) Calcul}

Une primitive de \(x^n\) est
\[
\frac{x^{n+1}}{n+1}.
\]
Donc
\[
I_n=\left[\frac{x^{n+1}}{n+1}\right]_0^1=\frac{1}{n+1}.
\]

\textbf{2) Étude de la suite}

\[
I_n=\frac{1}{n+1}.
\]
La suite est positive et décroissante.

\textbf{3) Limite}

\[
\lim_{n\to+\infty} I_n=\lim_{n\to+\infty}\frac{1}{n+1}=0.
\]

\[
\boxed{I_n=\frac{1}{n+1},\qquad \lim I_n=0.}
\]

% =========================================================
\begin{chapitrebox}{Chapitre 6 --- Équations différentielles}
\end{chapitrebox}

\subsection*{Exercice 1}
\[
y'+2y=0.
\]

C'est une équation différentielle linéaire du premier ordre à coefficients constants, de la forme
\[
y'+ay=0.
\]
La solution générale est
\[
y(x)=Ce^{-2x},\qquad C\in\R.
\]

\[
\boxed{y(x)=Ce^{-2x}.}
\]

% ---------------------------------------------------------
\subsection*{Exercice 2}
\[
y'-3y=0,\qquad y(0)=5.
\]

La solution générale est
\[
y(x)=Ce^{3x}.
\]
La condition initiale donne
\[
y(0)=C=5.
\]
Donc
\[
\boxed{y(x)=5e^{3x}.}
\]

% ---------------------------------------------------------
\subsection*{Exercice 3}
\[
y'+y=2.
\]

\begin{methodebox}
On cherche d'abord une solution particulière constante.
\end{methodebox}

Supposons \(y(x)=k\) constante. Alors \(y'=0\), donc
\[
0+k=2 \Rightarrow k=2.
\]
Une solution particulière est donc \(y_p(x)=2\).

La solution générale de l'équation homogène associée
\[
y'+y=0
\]
est
\[
y_h(x)=Ce^{-x}.
\]
Donc la solution générale est
\[
\boxed{y(x)=2+Ce^{-x}.}
\]

% ---------------------------------------------------------
\subsection*{Exercice 4}
\[
y'=ky.
\]

La solution générale est
\[
y(x)=Ce^{kx}.
\]
On demande les solutions croissantes.

On calcule
\[
y'(x)=kCe^{kx}.
\]
Comme \(e^{kx}>0\), le signe de \(y'\) est celui de \(kC\).

Donc :
\begin{itemize}
    \item si \(k>0\), les solutions croissantes sont celles pour lesquelles \(C>0\) ;
    \item si \(k<0\), les solutions croissantes sont celles pour lesquelles \(C<0\) ;
    \item si \(k=0\), toutes les solutions sont constantes, donc non strictement croissantes.
\end{itemize}

% ---------------------------------------------------------
\subsection*{Exercice 5}
\[
P'(t)=0{,}3P(t),\qquad P(0)=500.
\]

\textbf{1) Solution}

La solution générale est
\[
P(t)=Ce^{0{,}3t}.
\]
Avec \(P(0)=500\), on obtient
\[
C=500.
\]
Donc
\[
\boxed{P(t)=500e^{0{,}3t}.}
\]

\textbf{2) Quand dépasse-t-on 2000 ?}

On résout
\[
500e^{0{,}3t}>2000
\iff e^{0{,}3t}>4
\iff 0{,}3t>\ln 4
\iff t>\frac{\ln 4}{0{,}3}.
\]
Donc le dépassement a lieu pour
\[
\boxed{t>\dfrac{\ln 4}{0{,}3}\approx 4{,}62.}
\]

\textbf{3) Interprétation}

La population suit une croissance exponentielle continue de taux \(30\%\) par unité de temps.  
Le modèle suppose un environnement sans limitation de ressources.

% =========================================================
\begin{chapitrebox}{Chapitre 7 --- Probabilités}
\end{chapitrebox}

\subsection*{Exercice 1}
Deux dés équilibrés.

L'univers comporte \(6\times 6=36\) issues équiprobables.

\textbf{1) Somme égale à 7}

Les couples favorables sont :
\[
(1,6),(2,5),(3,4),(4,3),(5,2),(6,1).
\]
Il y en a \(6\). Donc
\[
P(S=7)=\frac{6}{36}=\frac16.
\]

\textbf{2) Somme \(\geq 10\)}

Somme 10 : \((4,6),(5,5),(6,4)\), soit 3 issues.  
Somme 11 : \((5,6),(6,5)\), soit 2 issues.  
Somme 12 : \((6,6)\), soit 1 issue.

Total : \(3+2+1=6\) issues. Donc
\[
P(S\geq 10)=\frac{6}{36}=\frac16.
\]

\textbf{3) Incompatibilité}

Un événement ne peut pas avoir une somme égale à 7 et en même temps une somme supérieure ou égale à 10.  
Donc les deux événements sont incompatibles.

% ---------------------------------------------------------
\subsection*{Exercice 2}
Urne : 5 rouges, 3 bleues, 2 vertes. Total : 10 boules.

\textbf{1) Rouge}

\[
P(R)=\frac{5}{10}=\frac12.
\]

\textbf{2) Non verte}

Il y a \(8\) boules non vertes. Donc
\[
P(\overline{V})=\frac{8}{10}=\frac45.
\]

\textbf{3) Rouge ou bleue}

Rouges ou bleues : \(5+3=8\) boules.
\[
P(R\cup B)=\frac{8}{10}=\frac45.
\]

% ---------------------------------------------------------
\subsection*{Exercice 3}
\[
P(A)=0{,}6,\quad P(B)=0{,}5,\quad P(A\cap B)=0{,}3.
\]

\textbf{1) Union}

\[
P(A\cup B)=P(A)+P(B)-P(A\cap B)=0{,}6+0{,}5-0{,}3=0{,}8.
\]

\textbf{2) Indépendance}

\(A\) et \(B\) sont indépendants si
\[
P(A\cap B)=P(A)P(B).
\]
Or
\[
P(A)P(B)=0{,}6\times 0{,}5=0{,}3=P(A\cap B).
\]
Donc \(A\) et \(B\) sont indépendants.

\textbf{3) Probabilité conditionnelle}

\[
P_A(B)=P(B\mid A)=\frac{P(A\cap B)}{P(A)}=\frac{0{,}3}{0{,}6}=0{,}5.
\]

% ---------------------------------------------------------
\subsection*{Exercice 4}
Maladie \(M\), test positif \(T\).

\[
P(M)=0{,}02,\qquad P(\overline M)=0{,}98,
\]
\[
P(T\mid M)=0{,}95,\qquad P(T\mid \overline M)=0{,}04.
\]

\textbf{1) Arbre pondéré}

On aurait :
\[
M \to T \text{ avec } 0{,}95,\qquad M \to \overline T \text{ avec } 0{,}05,
\]
\[
\overline M \to T \text{ avec } 0{,}04,\qquad \overline M \to \overline T \text{ avec } 0{,}96.
\]

\textbf{2) Probabilité d'être positif}

Par formule des probabilités totales :
\[
P(T)=P(M)P(T\mid M)+P(\overline M)P(T\mid \overline M).
\]
Donc
\[
P(T)=0{,}02\times 0{,}95+0{,}98\times 0{,}04=0{,}019+0{,}0392=0{,}0582.
\]

\textbf{3) Probabilité d'être malade sachant positif}

Par Bayes :
\[
P(M\mid T)=\frac{P(M\cap T)}{P(T)}
=\frac{P(M)P(T\mid M)}{P(T)}
=\frac{0{,}019}{0{,}0582}.
\]
Donc
\[
P(M\mid T)\approx 0{,}326.
\]

\[
\boxed{P(M\mid T)\approx 32{,}6\%.}
\]

\begin{logiquebox}
Un test peut être très sensible, et pourtant un test positif ne signifie pas forcément une très forte probabilité d'être malade, surtout si la maladie est rare.
\end{logiquebox}

% ---------------------------------------------------------
\subsection*{Exercice 5}
\[
P(A\cap B)=P(B)P_B(A).
\]

Par définition,
\[
P_B(A)=P(A\mid B)=\frac{P(A\cap B)}{P(B)} \quad \text{si } P(B)\neq 0.
\]
En multipliant par \(P(B)\), on obtient :
\[
\boxed{P(A\cap B)=P(B)P_B(A).}
\]

Pour une partition simple \(B,\overline B\),
\[
A=(A\cap B)\cup (A\cap \overline B),
\]
avec réunion disjointe, donc
\[
P(A)=P(A\cap B)+P(A\cap \overline B)
=P(B)P(A\mid B)+P(\overline B)P(A\mid \overline B).
\]

Exemple : si \(P(B)=0{,}4\), \(P(A\mid B)=0{,}7\), \(P(A\mid \overline B)=0{,}2\), alors
\[
P(A)=0{,}4\times 0{,}7+0{,}6\times 0{,}2=0{,}28+0{,}12=0{,}40.
\]

% =========================================================
\begin{chapitrebox}{Chapitre 8 --- Variables aléatoires et loi binomiale}
\end{chapitrebox}

\subsection*{Exercice 1}
On lance 5 fois une pièce équilibrée. \(X\) = nombre de piles.

Alors
\[
X\sim \mathcal{B}(5,\tfrac12).
\]

\textbf{1) Loi}

Pour \(k\in\{0,1,2,3,4,5\}\),
\[
P(X=k)=\binom{5}{k}\left(\frac12\right)^5.
\]

\textbf{2) \(P(X=3)\)}

\[
P(X=3)=\binom53\left(\frac12\right)^5=\frac{10}{32}=\frac{5}{16}.
\]

\textbf{3) \(P(X\geq 1)\)}

\[
P(X\geq 1)=1-P(X=0)=1-\left(\frac12\right)^5=1-\frac1{32}=\frac{31}{32}.
\]

% ---------------------------------------------------------
\subsection*{Exercice 2}
Défauts avec probabilité \(0{,}03\), 20 pièces, \(X\)= nombre de pièces défectueuses.

\textbf{1) Justification}

Chaque pièce peut être considérée comme :
\begin{itemize}
    \item un essai indépendant,
    \item à deux issues : défectueuse / non défectueuse,
    \item avec probabilité constante \(p=0{,}03\) de défaut.
\end{itemize}
Donc
\[
X\sim \mathcal{B}(20,0{,}03).
\]

\textbf{2) Paramètres}

\[
n=20,\qquad p=0{,}03.
\]

\textbf{3) \(P(X=0)\)}

\[
P(X=0)=(1-0{,}03)^{20}=0{,}97^{20}.
\]

\textbf{4) \(P(X\leq 2)\)}

\[
P(X\leq 2)=P(X=0)+P(X=1)+P(X=2).
\]
Soit
\[
P(X\leq 2)=0{,}97^{20}+\binom{20}{1}(0{,}03)(0{,}97)^{19}
+\binom{20}{2}(0{,}03)^2(0{,}97)^{18}.
\]

% ---------------------------------------------------------
\subsection*{Exercice 3}
\[
X\sim \mathcal{B}(n,p).
\]

\textbf{1) Espérance}

La formule de cours est
\[
\boxed{E(X)=np.}
\]

\textbf{2) Cas \(n=10\), \(p=0{,}2\)}

\[
E(X)=10\times 0{,}2=2.
\]

\textbf{3) Interprétation}

Sur un très grand nombre de répétitions de l'expérience, le nombre moyen de succès se rapproche de \(2\).

% ---------------------------------------------------------
\subsection*{Exercice 4}
Probabilité de victoire : \(0{,}4\), 8 parties, \(X\)= nombre de victoires.

\[
X\sim \mathcal{B}(8,0{,}4).
\]

\textbf{1) \(P(X=4)\)}

\[
P(X=4)=\binom84 (0{,}4)^4(0{,}6)^4.
\]

\textbf{2) \(P(X\geq 4)\)}

\[
P(X\geq 4)=\sum_{k=4}^{8}\binom{8}{k}(0{,}4)^k(0{,}6)^{8-k}.
\]

\textbf{3) Nombre moyen}

\[
E(X)=8\times 0{,}4=3{,}2.
\]

% ---------------------------------------------------------
\subsection*{Exercice 5}
\[
X\sim \mathcal{B}(6,p),\qquad P(X=0)=\frac1{64}.
\]

Or
\[
P(X=0)=(1-p)^6.
\]
Donc
\[
(1-p)^6=\frac1{64}=2^{-6}.
\]
Ainsi
\[
1-p=\frac12 \quad \text{(car } 1-p\in[0,1]\text{)},
\]
d'où
\[
p=\frac12.
\]

\textbf{2) \(P(X=2)\)}

\[
P(X=2)=\binom62 \left(\frac12\right)^2\left(\frac12\right)^4
=\binom62 \left(\frac12\right)^6
=\frac{15}{64}.
\]

\textbf{3) Espérance}

\[
E(X)=np=6\times \frac12=3.
\]

% =========================================================
\begin{chapitrebox}{Chapitre 9 --- Géométrie dans l'espace}
\end{chapitrebox}

\subsection*{Exercice 1}
\[
A(1,0,2),\quad B(2,1,3),\quad C(0,1,1).
\]

\textbf{1) Vecteurs}

\[
\overrightarrow{AB}=(2-1,\;1-0,\;3-2)=(1,1,1),
\]
\[
\overrightarrow{AC}=(0-1,\;1-0,\;1-2)=(-1,1,-1).
\]

\textbf{2) Alignement}

Les points \(A,B,C\) seraient alignés si \(\overrightarrow{AB}\) et \(\overrightarrow{AC}\) étaient colinéaires.

Or s'il existait \(\lambda\) tel que
\[
(-1,1,-1)=\lambda(1,1,1),
\]
alors il faudrait simultanément
\[
-1=\lambda,\qquad 1=\lambda,\qquad -1=\lambda,
\]
ce qui est impossible.

Donc les vecteurs ne sont pas colinéaires, ainsi
\[
\boxed{A,B,C \text{ ne sont pas alignés}.}
\]

% ---------------------------------------------------------
\subsection*{Exercice 2}
Droite passant par \(A(1,2,3)\) de vecteur directeur \((2,-1,4)\).

Une représentation paramétrique est :
\[
\boxed{
\begin{cases}
x=1+2t\\
y=2-t\\
z=3+4t
\end{cases}
\qquad t\in\R.
}
\]

% ---------------------------------------------------------
\subsection*{Exercice 3}
Plan
\[
(P): 2x-y+z-3=0.
\]

\textbf{1) Appartenance de \(A(1,0,1)\)}

On remplace :
\[
2\cdot1-0+1-3=2+1-3=0.
\]
Donc
\[
A\in (P).
\]

\textbf{2) Vecteur normal}

Dans l'équation
\[
ax+by+cz+d=0,
\]
un vecteur normal est \((a,b,c)\). Donc ici un vecteur normal est
\[
\boxed{(2,-1,1).}
\]

\textbf{3) Parallélisme avec une droite}

Une droite est parallèle à un plan si son vecteur directeur est orthogonal à un vecteur normal du plan.

Ici le vecteur directeur proposé est précisément \((2,-1,1)\), qui est parallèle au vecteur normal.  
Donc la droite est \textbf{perpendiculaire} au plan, et non parallèle.

% ---------------------------------------------------------
\subsection*{Exercice 4}
\[
d:\begin{cases}
x=1+t\\
y=2-t\\
z=3+2t
\end{cases}
\qquad
d':\begin{cases}
x=2+s\\
y=1+2s\\
z=5+s
\end{cases}
\]

\textbf{1) Sont-elles parallèles ?}

Vecteur directeur de \(d\) :
\[
\vec u=(1,-1,2).
\]
Vecteur directeur de \(d'\) :
\[
\vec v=(1,2,1).
\]
Ces vecteurs ne sont pas colinéaires, donc les droites ne sont pas parallèles.

\textbf{2) Sont-elles sécantes ?}

On cherche \(t,s\) tels que
\[
1+t=2+s,\qquad 2-t=1+2s,\qquad 3+2t=5+s.
\]
Soit
\[
t-s=1 \tag{1}
\]
\[
-t-2s=-1 \tag{2}
\]
\[
2t-s=2 \tag{3}
\]
De (1), \(t=1+s\).  
Dans (2) :
\[
-(1+s)-2s=-1
\iff -1-3s=-1
\iff s=0.
\]
Donc \(t=1\).  
Vérifions dans (3) :
\[
2\times 1-0=2,
\]
c'est vrai.

Donc les droites sont sécantes, au point obtenu pour \(t=1\) :
\[
(2,1,5).
\]

\[
\boxed{d \text{ et } d' \text{ sont sécantes en } (2,1,5).}
\]

% ---------------------------------------------------------
\subsection*{Exercice 5}
Dans un cube, montrer qu'une droite est perpendiculaire à un plan.

\begin{methodebox}
Le critère fondamental est :
\[
\text{une droite est perpendiculaire à un plan si elle est perpendiculaire à deux droites sécantes de ce plan.}
\]
\end{methodebox}

Dans un cube \(ABCDEFGH\), par exemple, pour montrer que la droite \(AE\) est perpendiculaire au plan \((ABCD)\), on remarque que :
\begin{itemize}
    \item \(AE\perp AB\),
    \item \(AE\perp AD\),
\end{itemize}
or \(AB\) et \(AD\) sont deux droites sécantes du plan \((ABCD)\).  
Donc
\[
\boxed{AE\perp (ABCD).}
\]

% =========================================================
\begin{chapitrebox}{Chapitre 10 --- Nombres complexes}
\end{chapitrebox}

\subsection*{Exercice 1}
\[
z^2+4=0.
\]

\[
z^2=-4=4i^2.
\]
Donc
\[
z=\pm 2i.
\]

\[
\boxed{z=2i \quad \text{ou} \quad z=-2i.}
\]

% ---------------------------------------------------------
\subsection*{Exercice 2}
\[
(2+3i)(1-4i).
\]

Développons :
\[
(2+3i)(1-4i)=2-8i+3i-12i^2.
\]
Or \(i^2=-1\), donc
\[
-12i^2=12.
\]
Ainsi
\[
2+12-5i=14-5i.
\]

\[
\boxed{(2+3i)(1-4i)=14-5i.}
\]

% ---------------------------------------------------------
\subsection*{Exercice 3}
\[
z^2-2z+5=0.
\]

Discriminant :
\[
\Delta=(-2)^2-4\times 1\times 5=4-20=-16.
\]
Donc
\[
z=\frac{2\pm \sqrt{-16}}{2}=\frac{2\pm 4i}{2}=1\pm 2i.
\]

\[
\boxed{z=1+2i \quad \text{ou} \quad z=1-2i.}
\]

% ---------------------------------------------------------
\subsection*{Exercice 4}
\[
z=1+i\sqrt3.
\]

\textbf{1) Module}

\[
|z|=\sqrt{1^2+(\sqrt3)^2}=\sqrt{4}=2.
\]

\textbf{2) Argument}

\[
\cos \theta=\frac12,\qquad \sin\theta=\frac{\sqrt3}{2}.
\]
Donc un argument est
\[
\theta=\frac{\pi}{3}.
\]

\textbf{3) Forme trigonométrique}

\[
\boxed{z=2\left(\cos\frac{\pi}{3}+i\sin\frac{\pi}{3}\right).}
\]

% ---------------------------------------------------------
\subsection*{Exercice 5}
Affixes :
\[
A(1),\quad B(i),\quad C(1+i).
\]

Dans le plan, cela donne :
\[
A(1,0),\qquad B(0,1),\qquad C(1,1).
\]

\textbf{1) Placement}

Les points sont donc faciles à placer dans un repère orthonormé.

\textbf{2) Triangle rectangle}

Calculons :
\[
\overrightarrow{CA}=A-C=(0,-1),\qquad \overrightarrow{CB}=B-C=(-1,0).
\]
Le produit scalaire vaut
\[
\overrightarrow{CA}\cdot \overrightarrow{CB}=0\times (-1)+(-1)\times 0=0.
\]
Donc
\[
\boxed{ABC \text{ est rectangle en } C.}
\]

\textbf{3) Isocèle}

\[
CA=\sqrt{0^2+(-1)^2}=1,\qquad CB=\sqrt{(-1)^2+0^2}=1.
\]
Donc
\[
CA=CB.
\]
Ainsi
\[
\boxed{ABC \text{ est isocèle en } C.}
\]

% =========================================================
\begin{chapitrebox}{Chapitre 11 --- Algorithmique et raisonnement}
\end{chapitrebox}

\subsection*{Exercice 1}
Algorithme :
\[
u\leftarrow 1,\qquad u\leftarrow 2u+1 \text{ répété } n \text{ fois.}
\]

\textbf{1) Premières valeurs}

\[
u_0=1,
\]
\[
u_1=2\cdot 1+1=3,
\]
\[
u_2=2\cdot 3+1=7,
\]
\[
u_3=2\cdot 7+1=15.
\]
On reconnaît
\[
1,3,7,15,\dots
\]

\textbf{2) Conjecture}

On conjecture
\[
u_n=2^{n+1}-1.
\]

\textbf{3) Démonstration par récurrence}

\emph{Initialisation :}
\[
u_0=1=2^{1}-1.
\]
La formule est vraie au rang \(0\).

\emph{Hérédité :} supposons
\[
u_n=2^{n+1}-1.
\]
Alors
\[
u_{n+1}=2u_n+1=2(2^{n+1}-1)+1=2^{n+2}-2+1=2^{n+2}-1.
\]
La propriété est vraie au rang \(n+1\).

Donc, par récurrence,
\[
\boxed{u_n=2^{n+1}-1.}
\]

% ---------------------------------------------------------
\subsection*{Exercice 2}
Trouver le plus petit \(n\) tel que
\[
\frac{n}{2^n}<10^{-4}.
\]

Un algorithme possible :

\begin{verbatim}
n <- 0
tant que n / 2^n >= 10^(-4) faire
    n <- n + 1
fin tant que
afficher n
\end{verbatim}

\begin{logiquebox}
L'idée clé d'un tel algorithme est de partir d'une valeur simple, puis d'augmenter \(n\) tant que la condition recherchée n'est pas satisfaite.
\end{logiquebox}

% ---------------------------------------------------------
\subsection*{Exercice 3}
Suite définie par récurrence.

\textbf{Calcul de \(u_n\)}

\begin{verbatim}
u <- u0
pour k allant de 1 à n faire
    u <- expression_de_u_en_fonction_de_u
fin pour
afficher u
\end{verbatim}

\textbf{Premier rang tel que \(u_n>100\)}

\begin{verbatim}
u <- u0
n <- 0
tant que u <= 100 faire
    u <- expression_de_u_en_fonction_de_u
    n <- n + 1
fin tant que
afficher n
\end{verbatim}

% ---------------------------------------------------------
\subsection*{Exercice 4}
Approcher la solution de
\[
x^3+x-1=0
\quad \text{sur } [0,1].
\]

Posons
\[
f(x)=x^3+x-1.
\]

\textbf{1) Existence et unicité}

\(f\) est continue sur \([0,1]\).  
On a
\[
f(0)=-1<0,\qquad f(1)=1>0.
\]
Par théorème des valeurs intermédiaires, il existe au moins une solution sur \([0,1]\).

De plus
\[
f'(x)=3x^2+1>0.
\]
Donc \(f\) est strictement croissante sur \([0,1]\), d'où l'unicité de la solution.

\textbf{2) Dichotomie}

On coupe l'intervalle en deux, on regarde le signe au milieu, puis on garde la moitié où le changement de signe subsiste.

\textbf{3) Algorithme}

\begin{verbatim}
a <- 0
b <- 1
tant que b - a > precision faire
    m <- (a + b)/2
    si f(a) * f(m) <= 0 alors
        b <- m
    sinon
        a <- m
    fin si
fin tant que
afficher (a+b)/2
\end{verbatim}

% ---------------------------------------------------------
\subsection*{Exercice 5}
\(f\) strictement croissante sur \([a,b]\).

\textbf{1) Pourquoi au plus une solution ?}

Supposons qu'il existe deux solutions distinctes \(x_1<x_2\) de \(f(x)=0\).  
Comme \(f\) est strictement croissante,
\[
f(x_1)<f(x_2).
\]
Mais
\[
f(x_1)=0=f(x_2),
\]
contradiction. Donc il y a au plus une solution.

\textbf{2) Méthode d'approximation si \(f(a)<0<f(b)\)}

Comme \(f\) est continue et strictement croissante, il existe une unique solution.  
On peut utiliser la dichotomie sur \([a,b]\).

% =========================================================
\begin{chapitrebox}{Chapitre 12 --- Arithmétique}
\end{chapitrebox}

\subsection*{Exercice 1}
Reste de \(2^{10}\) dans la division par \(7\).

Calculons modulo \(7\) :
\[
2^1=2,\quad 2^2=4,\quad 2^3=8\equiv 1 \pmod 7.
\]
Donc
\[
2^{10}=2^{9}2=(2^3)^3\cdot 2 \equiv 1^3\cdot 2\equiv 2 \pmod 7.
\]

\[
\boxed{\text{Le reste est } 2.}
\]

% ---------------------------------------------------------
\subsection*{Exercice 2}
Si \(n\) est impair, montrer que \(n^2\) est impair.

Si \(n\) est impair, il existe \(k\in\Z\) tel que
\[
n=2k+1.
\]
Alors
\[
n^2=(2k+1)^2=4k^2+4k+1=2(2k^2+2k)+1.
\]
Donc \(n^2\) est de la forme \(2m+1\), donc impair.

\[
\boxed{n \text{ impair } \Rightarrow n^2 \text{ impair}.}
\]

% ---------------------------------------------------------
\subsection*{Exercice 3}
Montrer que
\[
n(n+1)
\]
est pair.

Parmi deux entiers consécutifs \(n\) et \(n+1\), l'un est nécessairement pair.  
Le produit d'un entier pair par n'importe quel entier est pair.

Donc
\[
\boxed{n(n+1)\text{ est pair pour tout } n\in\Z.}
\]

% ---------------------------------------------------------
\subsection*{Exercice 4}
Déterminer les entiers \(n\) tels que
\[
n^2-n
\]
soit divisible par \(2\).

On factorise :
\[
n^2-n=n(n-1).
\]
Or \(n\) et \(n-1\) sont deux entiers consécutifs, donc l'un des deux est pair.  
Leur produit est donc pair.

Ainsi, pour tout entier \(n\),
\[
\boxed{n^2-n \text{ est divisible par } 2.}
\]

% ---------------------------------------------------------
\subsection*{Exercice 5}
Montrer que
\[
n^3-n
\]
est divisible par \(6\).

On factorise :
\[
n^3-n=n(n^2-1)=n(n-1)(n+1).
\]
C'est le produit de trois entiers consécutifs.

\textbf{Divisibilité par 2 :}  
Parmi trois entiers consécutifs, il y en a au moins un qui est pair. Donc le produit est divisible par \(2\).

\textbf{Divisibilité par 3 :}  
Parmi trois entiers consécutifs, il y en a exactement un qui est multiple de \(3\). Donc le produit est divisible par \(3\).

Comme \(2\) et \(3\) sont premiers entre eux, le produit est divisible par \(6\).

Ainsi,
\[
\boxed{6 \mid (n^3-n)\quad \text{pour tout entier } n.}
\]

% =========================================================
\section*{Conclusion générale}

Ce corrigé met en évidence plusieurs idées essentielles :

\begin{itemize}
    \item \textbf{Toujours vérifier le domaine de définition} avant de résoudre.
    \item \textbf{Ne pas confondre expression simplifiée et fonction initiale}.
    \item \textbf{Pour une suite récurrente}, penser à :
    \begin{itemize}
        \item montrer qu'elle est bien définie,
        \item chercher un intervalle stable,
        \item étudier la monotonie,
        \item seulement ensuite conclure sur la limite.
    \end{itemize}
    \item \textbf{Pour une équation}, distinguer :
    \begin{itemize}
        \item existence,
        \item unicité,
        \item calcul effectif de la solution.
    \end{itemize}
    \item \textbf{En probabilités}, bien identifier :
    \begin{itemize}
        \item l'univers,
        \item les événements,
        \item les formules adaptées.
    \end{itemize}
    \item \textbf{En arithmétique}, factoriser est souvent la bonne idée de départ.
\end{itemize}

\vspace{0.5em}

\begin{methodebox}
Pour aller encore plus loin, une excellente habitude de rédaction consiste à toujours répondre à ces trois questions :
\begin{enumerate}
    \item \textbf{Que cherche-t-on ?}
    \item \textbf{Quel outil a-t-on le droit d'utiliser ?}
    \item \textbf{Pourquoi cet outil est-il pertinent ici ?}
\end{enumerate}
\end{methodebox}

\end{document}